% ======================================================================
% 8. DETECTION  (budget: 1.5 pages)
% Authority: every event ID mentioned must be confirmed against Microsoft docs.
% ======================================================================
\section{Detection}
\label{sec:detection}

\subsection{Host telemetry that matters}
% TODO
% - 4624 logon type 3/9 with NTLM as the authentication package, and the
%   detail fields worth alerting on (workstation name vs source IP mismatch,
%   logon without a preceding interactive session).
% - 4672 special privileges assigned, following a type-3 logon.
% - 4776 credential validation on the DC and how it differs from 4624.
% - 4768/4769 for the PtK side, including cipher/etype fields.
% - 4648 explicit credential use as a baseline comparator.

\subsection{What the attack cannot hide}
% TODO
% - The target always logs a successful network logon; the DC always validates.
%   The attacker controls neither. This is the structural weakness of PtH.
% - Therefore detection is about baselining "who logs on where with NTLM",
%   not about catching a signature.

\subsection{Network-level signals}
% TODO
% - NTLMSSP in clear where the carrier is not encrypted; Negotiate downgrade.
% - NTLM auditing on the DC as the cheapest visibility win.

\subsection{Behavioural and deception controls}
% TODO
% - Honeytoken accounts / canary credentials in LSASS; any use is high signal.
% - Baseline deviation: NTLM logon to a server that normally only does Kerberos.
% - Short paragraph on EDR: what a memory read of LSASS looks like.
